\documentclass[11pt]{article}
\usepackage[margin=1in]{geometry}
\usepackage{enumitem}
\usepackage{hyperref}
\usepackage{graphicx}
\usepackage{array}
\usepackage{multicol}
\usepackage{longtable}
\usepackage{titlesec}
\usepackage{listings}
\usepackage{xcolor}
\usepackage[a4paper, margin=1in]{geometry}
\usepackage{float}      % For [H] placement
\usepackage{caption}
\usepackage{subcaption}
\usepackage{amsmath}
\usepackage{setspace}
\usepackage[margin=1in]{geometry}
\usepackage{enumitem}
\usepackage{hyperref}
\usepackage{graphicx}
\usepackage{array}
\usepackage{multicol}
\usepackage{longtable}
\usepackage{titlesec}
\usepackage{array} % add this in preamble

% Custom colors for code
\definecolor{codegreen}{rgb}{0,0.6,0}
\definecolor{codegray}{rgb}{0.5,0.5,0.5}
\definecolor{codepurple}{rgb}{0.58,0,0.82}
\definecolor{backcolour}{rgb}{0.95,0.95,0.92}

\lstset{
    backgroundcolor=\color{backcolour},   
    commentstyle=\color{codegreen},
    keywordstyle=\color{magenta},
    numberstyle=\tiny\color{codegray},
    stringstyle=\color{codepurple},
    basicstyle=\ttfamily\footnotesize,
    breakatwhitespace=false,         
    breaklines=true,                 
    captionpos=b,                    
    keepspaces=true,                 
    numbers=none,                    
    numbersep=5pt,                  
    showspaces=false,                
    showstringspaces=false,
    showtabs=false,                  
    tabsize=2,
    language=Python
}

\begin{document}

%==================================================
\begin{center}
    \large \textbf{Sri Sivasubramaniya Nadar College of Engineering, Chennai} \\
    (An Autonomous Institution Affiliated to Anna University) \\
    \vspace{0.3cm}
\end{center}

\begin{table}[!h]
\renewcommand{\arraystretch}{1.5}
\resizebox{\textwidth}{!}{%
\begin{tabular}{|l|cll|}
\hline
Degree \& Branch & \multicolumn{1}{c|}{B.E. Computer Science \& Engineering} & Semester & VI \\ \hline
Subject Code \& Name & \multicolumn{3}{c|}{UCS2612 -- Machine Learning Algorithms Laboratory} \\ \hline
Academic Year & 2025--2026 (Even) & Batch & 2023--2027 \\ \hline
Due Date & \multicolumn{3}{c|}{29-03-2026} \\ \hline
Name & \multicolumn{3}{c|}{Priya Verma - 3122235001102} \\ \hline
\end{tabular}
}
\end{table}

\begin{center}
    \textbf{\Large Experiment 8}
\end{center}

\begin{center}
    \textbf{Effect of PCA on Regression and Classification Models}
\end{center}

%==================================================
\section{Objective}
To study the impact of \textbf{Principal Component Analysis (PCA)} 
on the performance of machine learning models.

Students must:

\begin{enumerate}[label=\alph*)]
    \item Train models in the original feature space (Without PCA)
    \item Train models in reduced feature space (With PCA)
    \item Compare performance using 5-fold cross-validation
\end{enumerate}

%==================================================
\section{Dataset}
Dataset: \textbf{PCA\_Full\_Lab\_Dataset\_1000\_Samples.csv}

\textbf{Features:} Academic and performance-related attributes  
\textbf{Regression Target:} Final\_Score\_Regression  
\textbf{Classification Target:} Performance\_Level\_Classification  

%==================================================
\section{Implementation Details}
The following code snippets illustrate the core implementation of the classifiers:

\subsection*{Important Libraries to be included}
\begin{lstlisting}[language=Python]
import pandas as pd
import numpy as np
import matplotlib.pyplot as plt

from sklearn.model_selection import cross_val_score, KFold, GridSearchCV
from sklearn.preprocessing import StandardScaler, LabelEncoder
from sklearn.decomposition import PCA

from sklearn.linear_model import LinearRegression, LogisticRegression
from sklearn.ensemble import RandomForestRegressor
from sklearn.svm import SVC

from sklearn.metrics import make_scorer, mean_squared_error, r2_score, accuracy_score, f1_score

plt.rcParams.update({
    "font.family": "Arial",
    "font.size": 15,
    "axes.titlesize": 15,
    "axes.titleweight": "bold",
    "axes.labelsize": 15,
    "axes.labelweight": "bold",
    "xtick.labelsize": 13,
    "ytick.labelsize": 13
})
\end{lstlisting}

\subsection*{1. Load the dataset}
\begin{lstlisting}[language=Python]
# Load dataset
df = pd.read_csv("Dataset.csv")

print("Dataset Shape:", df.shape)
df.head()
# Number of samples and features
print("Number of samples:", df.shape[0])
print("Number of features:", df.shape[1]-2)

# Target distribution
print("\nClassification Target Distribution")
print(df['Performance_Level_Classification'].value_counts())
\end{lstlisting}

\subsection*{2. Preprocessing Steps}
\begin{lstlisting}[language=Python]
# 2a. Identify target columns 
# Adjust these column names if your CSV uses different names
REGRESSION_TARGET    = "Final Score"        # continuous target
CLASSIFICATION_TARGET = "Performance Level" # categorical target

# Auto-detect if column names differ
all_cols = df.columns.tolist()
if REGRESSION_TARGET not in all_cols:
    # Try to find a numeric column that could be regression target
    num_cols = df.select_dtypes(include=[np.number]).columns.tolist()
    REGRESSION_TARGET = num_cols[-1]
    print(f"[Auto-detected] Regression target    : '{REGRESSION_TARGET}'")
 
if CLASSIFICATION_TARGET not in all_cols:
    # Try to find a categorical/object column
    cat_cols = df.select_dtypes(include=["object", "category"]).columns.tolist()
    if cat_cols:
        CLASSIFICATION_TARGET = cat_cols[-1]
    else:
        # Fall back to binning the regression target
        CLASSIFICATION_TARGET = "Performance_Level_Binned"
        df[CLASSIFICATION_TARGET] = pd.qcut(
            df[REGRESSION_TARGET], q=3, labels=["Low", "Medium", "High"]
        )
    print(f"[Auto-detected] Classification target: '{CLASSIFICATION_TARGET}'")
 
print(f"\nRegression    target : '{REGRESSION_TARGET}'")
print(f"Classification target: '{CLASSIFICATION_TARGET}'")

# 2b. Separate feature columns 
drop_cols = [REGRESSION_TARGET, CLASSIFICATION_TARGET]
drop_cols = [c for c in drop_cols if c in df.columns]
X_raw = df.drop(columns=drop_cols)
 
y_reg  = df[REGRESSION_TARGET].copy()
y_clf  = df[CLASSIFICATION_TARGET].copy()

# 2c. Handle missing values 
missing_before = X_raw.isnull().sum().sum()
print(f"\nMissing values before imputation : {missing_before}")

# Numeric columns → median imputation
num_feat = X_raw.select_dtypes(include=[np.number]).columns
X_raw[num_feat] = X_raw[num_feat].fillna(X_raw[num_feat].median())

# Categorical columns → mode imputation then encode
cat_feat = X_raw.select_dtypes(include=["object", "category"]).columns
for col in cat_feat:
    X_raw[col] = X_raw[col].fillna(X_raw[col].mode()[0])
    X_raw[col] = LabelEncoder().fit_transform(X_raw[col].astype(str))
 
y_reg = y_reg.fillna(y_reg.median())
 
missing_after = X_raw.isnull().sum().sum()
print(f"Missing values after  imputation : {missing_after}")

# 2d. Encode classification target 
le = LabelEncoder()
y_clf_enc = le.fit_transform(y_clf.astype(str))
print(f"\nClassification classes : {list(le.classes_)}")

# 2e. Dataset statistics 
print(f"\nFeature matrix shape  : {X_raw.shape}")
print(f"\nRegression target distribution:")
print(f"  Mean={y_reg.mean():.2f}, Std={y_reg.std():.2f}, "
      f"Min={y_reg.min():.2f}, Max={y_reg.max():.2f}")
print(f"\nClassification target distribution:")
for cls, cnt in zip(le.classes_, np.bincount(y_clf_enc)):
    print(f"  {cls}: {cnt} samples ({100*cnt/len(y_clf_enc):.1f}%)")

# 2f. Standardize features
scaler = StandardScaler()
X_scaled = scaler.fit_transform(X_raw)
X_scaled = pd.DataFrame(X_scaled, columns=X_raw.columns)
print(f"\nStandardization applied — mean≈{X_scaled.values.mean():.4f}, "
      f"std≈{X_scaled.values.std():.4f}")
\end{lstlisting}

\subsection*{3. PCA Implementation}
\begin{lstlisting}[language=Python]
# Full PCA to study explained variance
pca_full = PCA(random_state=42)
pca_full.fit(X_scaled)
 
explained_var       = pca_full.explained_variance_ratio_
cumulative_var      = np.cumsum(explained_var)
n_components_95     = np.argmax(cumulative_var >= 0.95) + 1
n_features          = X_scaled.shape[1]
 
print(f"Total features               : {n_features}")
print(f"Components for 95% variance  : {n_components_95}")
print(f"Explained variance at 95% PC : {cumulative_var[n_components_95-1]*100:.2f}%")

# 3a. Scree Plot 
fig, axes = plt.subplots(1, 2, figsize=(14, 5))
 
axes[0].bar(range(1, min(21, n_features+1)),
            explained_var[:20] * 100,
            color="steelblue", edgecolor="black", alpha=0.8)
axes[0].set_title("Scree Plot: Individual Explained Variance", fontweight="bold")
axes[0].set_xlabel("Principal Component")
axes[0].set_ylabel("Explained Variance (%)")
axes[0].grid(axis="y", alpha=0.4)
 
axes[1].plot(range(1, n_features+1), cumulative_var * 100,
             "bo-", markersize=4, linewidth=2)
axes[1].axhline(y=95, color="red", linestyle="--", linewidth=1.5, label="95% threshold")
axes[1].axvline(x=n_components_95, color="green", linestyle="--",
                linewidth=1.5, label=f"k={n_components_95} components")
axes[1].fill_between(range(1, n_features+1), cumulative_var * 100,
                     alpha=0.15, color="steelblue")
axes[1].set_title("Cumulative Explained Variance", fontweight="bold")
axes[1].set_xlabel("Number of Components")
axes[1].set_ylabel("Cumulative Variance (%)")
axes[1].legend()
axes[1].grid(alpha=0.4)
 
plt.suptitle("PCA Scree Plot Analysis", fontsize=14, fontweight="bold")
plt.tight_layout()
plt.savefig("pca_scree_plot.png", dpi=150)
plt.show()
print("Saved: pca_scree_plot.png")

# 3b. Apply PCA (retain 95% variance) 
pca = PCA(n_components=n_components_95, random_state=42)
X_pca = pca.fit_transform(X_scaled)
X_pca = pd.DataFrame(X_pca, columns=[f"PC{i+1}" for i in range(n_components_95)])
 
print(f"\nTable 1: PCA Summary")
print(f"{'Components Chosen':<22} {n_components_95}")
print(f"{'Explained Variance (%)':<22} {cumulative_var[n_components_95-1]*100:.2f}%")
print(f"{'Justification':<22} Retain 95% of total variance with minimal information loss.")
 
X_no_pca = X_scaled   # Original scaled features
X_with_pca = X_pca    # PCA-reduced features

kf = KFold(n_splits=5, shuffle=True, random_state=42)

def cv_regression(model, X, y, kf):
    """Return per-fold MSE and R2."""
    mse_folds, r2_folds = [], []
    for train_idx, test_idx in kf.split(X):
        X_tr, X_te = X.iloc[train_idx], X.iloc[test_idx]
        y_tr, y_te = y.iloc[train_idx], y.iloc[test_idx]
        model.fit(X_tr, y_tr)
        y_pred = model.predict(X_te)
        mse_folds.append(mean_squared_error(y_te, y_pred))
        r2_folds.append(r2_score(y_te, y_pred))
    return np.array(mse_folds), np.array(r2_folds)

 def cv_classification(model, X, y, kf):
    """Return per-fold Accuracy and F1."""
    acc_folds, f1_folds = [], []
    for train_idx, test_idx in kf.split(X):
        X_tr, X_te = X.iloc[train_idx], X.iloc[test_idx]
        y_tr, y_te = y[train_idx], y[test_idx]
        model.fit(X_tr, y_tr)
        y_pred = model.predict(X_te)
        acc_folds.append(accuracy_score(y_te, y_pred))
        f1_folds.append(f1_score(y_te, y_pred, average="weighted", zero_division=0))
    return np.array(acc_folds), np.array(f1_folds)
 
y_reg_s  = pd.Series(y_reg.values)
y_clf_s  = np.array(y_clf_enc)
\end{lstlisting}

\subsection*{4. Regression Models}
\begin{lstlisting}[language=Python]
# 5a. Linear Regression 
print("\n--- Linear Regression ---")
lr = LinearRegression()
 
lr_mse_no,  lr_r2_no  = cv_regression(lr, X_no_pca,   y_reg_s, kf)
lr_mse_pca, lr_r2_pca = cv_regression(lr, X_with_pca, y_reg_s, kf)
 
print(f"  No PCA  — MSE: {lr_mse_no.mean():.4f} ± {lr_mse_no.std():.4f}  "
      f"R²: {lr_r2_no.mean():.4f} ± {lr_r2_no.std():.4f}")
print(f"  With PCA— MSE: {lr_mse_pca.mean():.4f} ± {lr_mse_pca.std():.4f}  "
      f"R²: {lr_r2_pca.mean():.4f} ± {lr_r2_pca.std():.4f}")

# 5b. Random Forest Regressor 
print("\n--- Random Forest Regressor (n_estimators=100, max_depth=10) ---")
rf = RandomForestRegressor(n_estimators=100, max_depth=10, random_state=42, n_jobs=-1)
 
rf_mse_no,  rf_r2_no  = cv_regression(rf, X_no_pca,   y_reg_s, kf)
rf_mse_pca, rf_r2_pca = cv_regression(rf, X_with_pca, y_reg_s, kf)
 
print(f"  No PCA  — MSE: {rf_mse_no.mean():.4f} ± {rf_mse_no.std():.4f}  "
      f"R²: {rf_r2_no.mean():.4f} ± {rf_r2_no.std():.4f}")
print(f"  With PCA— MSE: {rf_mse_pca.mean():.4f} ± {rf_mse_pca.std():.4f}  "
      f"R²: {rf_r2_pca.mean():.4f} ± {rf_r2_pca.std():.4f}")

# Table 2 
print("\n\nTable 2: 5-Fold CV Results — Regression")
print(f"{'Model':<22} {'Metric':<8} {'Fold Avg (No PCA)':>20} {'Fold Avg (With PCA)':>22} Observation")
print("-" * 100)
 
reg_rows = [
    ("Linear Regression", "MSE", lr_mse_no.mean(), lr_mse_pca.mean(),
     "PCA may reduce MSE by removing multicollinear features"),
    ("Linear Regression", "R²",  lr_r2_no.mean(),  lr_r2_pca.mean(),
     "R² typically stable or improved with PCA"),
    ("Random Forest",     "MSE", rf_mse_no.mean(), rf_mse_pca.mean(),
     "RF less sensitive to PCA; may slightly worsen"),
    ("Random Forest",     "R²",  rf_r2_no.mean(),  rf_r2_pca.mean(),
     "Tree models handle raw features well"),
]
for row in reg_rows:
    model, metric, no_pca, with_pca, obs = row
    print(f"{model:<22} {metric:<8} {no_pca:>20.4f} {with_pca:>22.4f}  {obs}")

# Per-fold detail
print("\nPer-Fold MSE — Linear Regression:")
for i, (a, b) in enumerate(zip(lr_mse_no, lr_mse_pca), 1):
    print(f"  Fold {i}: No PCA={a:.4f}  With PCA={b:.4f}  Diff={b-a:+.4f}")
 
print("\nPer-Fold R² — Random Forest:")
for i, (a, b) in enumerate(zip(rf_r2_no, rf_r2_pca), 1):
    print(f"  Fold {i}: No PCA={a:.4f}  With PCA={b:.4f}  Diff={b-a:+.4f}")
\end{lstlisting}

\subsection*{5. Classification Models}
\begin{lstlisting}[language=Python]
X_no_pca_clf   = X_no_pca.copy()
X_with_pca_clf = X_with_pca.copy()
 
# 6a. Logistic Regression (C=1.0) 
print("\n--- Logistic Regression (C=1.0, max_iter=1000) ---")
log_reg = LogisticRegression(C=1.0, max_iter=1000, random_state=42)
 
lr_acc_no,  lr_f1_no  = cv_classification(log_reg, X_no_pca_clf,   y_clf_s, kf)
lr_acc_pca, lr_f1_pca = cv_classification(log_reg, X_with_pca_clf, y_clf_s, kf)
 
print(f"  No PCA  — Acc: {lr_acc_no.mean():.4f} ± {lr_acc_no.std():.4f}  "
      f"F1: {lr_f1_no.mean():.4f} ± {lr_f1_no.std():.4f}")
print(f"  With PCA— Acc: {lr_acc_pca.mean():.4f} ± {lr_acc_pca.std():.4f}  "
      f"F1: {lr_f1_pca.mean():.4f} ± {lr_f1_pca.std():.4f}")

# 6b. SVM (RBF kernel) 
print("\n--- SVM (kernel=RBF, C=1.0, gamma=scale) ---")
svm = SVC(kernel="rbf", C=1.0, gamma="scale", random_state=42)
 
svm_acc_no,  svm_f1_no  = cv_classification(svm, X_no_pca_clf,   y_clf_s, kf)
svm_acc_pca, svm_f1_pca = cv_classification(svm, X_with_pca_clf, y_clf_s, kf)
 
print(f"  No PCA  — Acc: {svm_acc_no.mean():.4f} ± {svm_acc_no.std():.4f}  "
      f"F1: {svm_f1_no.mean():.4f} ± {svm_f1_no.std():.4f}")
print(f"  With PCA— Acc: {svm_acc_pca.mean():.4f} ± {svm_acc_pca.std():.4f}  "
      f"F1: {svm_f1_pca.mean():.4f} ± {svm_f1_pca.std():.4f}")

# Table 3
print("\n\nTable 3: 5-Fold CV Results — Classification")
print(f"{'Model':<22} {'Metric':<10} {'Fold Avg (No PCA)':>20} {'Fold Avg (With PCA)':>22} Observation")
print("-" * 105)
 
clf_rows = [
    ("Logistic Regression", "Accuracy", lr_acc_no.mean(), lr_acc_pca.mean(),
     "Linear model benefits noticeably from PCA"),
    ("Logistic Regression", "F1-score", lr_f1_no.mean(),  lr_f1_pca.mean(),
     "PCA removes redundant features improving F1"),
    ("SVM",                 "Accuracy", svm_acc_no.mean(), svm_acc_pca.mean(),
     "SVM with RBF can handle high-dim, PCA may help speed"),
    ("SVM",                 "F1-score", svm_f1_no.mean(),  svm_f1_pca.mean(),
     "Kernel SVM partially benefits from dimensionality reduction"),
]
for row in clf_rows:
    model, metric, no_pca, with_pca, obs = row
    print(f"{model:<22} {metric:<10} {no_pca:>20.4f} {with_pca:>22.4f}  {obs}")
\end{lstlisting}


%==================================================
% -------------------- SECTION --------------------
\section{Visualizations}

% -------------------- FIGURES --------------------
\begin{figure}[H]
    \centering
    \includegraphics[width=\textwidth]{pca_scree_plot.png}
    \caption{PCA Scree Plot Analysis}
\end{figure}

\begin{figure}[H]
    \centering
    \includegraphics[width=\textwidth]{regression_comparison.png}
    \caption{Regression Results — No PCA vs With PCA}
\end{figure}

\begin{figure}[H]
    \centering
    \includegraphics[width=\textwidth]{classification_comparison.png}
    \caption{Classification Results — No PCA vs With PCA}
\end{figure}

\begin{figure}[H]
    \centering
    \includegraphics[width=\textwidth]{per_fold_comparison.png}
    \caption{Per-Fold Performance: No PCA vs With PCA}
\end{figure}

\begin{figure}[H]
    \centering
    \includegraphics[width=\textwidth]{variance_comparison.png}
    \caption{Std Dev Across Folds (Lower = More Stable)}
\end{figure}

\begin{figure}[H]
    \centering
    \includegraphics[width=\textwidth]{pca_2d_scatter.png}
    \caption{PCA 2D Projection (PC1 vs PC2)}
\end{figure}

\begin{figure}[H]
    \centering
    \includegraphics[width=\textwidth]{feature_correlation_heatmap.png}
    \caption{Feature Correlation Matrix (before PCA)}
\end{figure}

\begin{figure}[H]
    \centering
    \includegraphics[width=\textwidth]{target_distributions.png}
    \caption{Target Variable Distributions}
\end{figure}

%==================================================
\section{Performance Tables}
%==================================================

\begin{table}[H]
\centering
\renewcommand{\arraystretch}{1.3}
\caption{PCA Summary}
\begin{tabular}{|c|c|p{5cm}|}
\hline
Components Chosen & Explained Variance (\%) & Justification \\
\hline
7 & 95.90\% & Retain 95\% of total variance with minimal information loss. \\
\hline
\end{tabular}
\end{table}

%==================================================
\begin{table}[H]
\centering
\renewcommand{\arraystretch}{1.3}
\caption{5-Fold CV Results – Regression}
\begin{tabular}{|l|c|c|c|p{5cm}|}
\hline
Model & Metric & Fold Avg (No-PCA) & Fold Avg (With-PCA) & Observation \\
\hline
Linear Regression & MSE & 26.2338 & 26.2251 & PCA may reduce MSE by removing multicollinear features \\
Linear Regression & $R^2$ & 0.7479 & 0.7481 & R² typically stable or improved with PCA \\
Random Forest & MSE & 28.9206 & 28.7573 & RF less sensitive to PCA; may slightly worsen \\
Random Forest & $R^2$ & 0.7225 & 0.7241 & Tree models handle raw features well \\
\hline
\end{tabular}
\end{table}

%==================================================
\begin{table}[H]
\centering
\renewcommand{\arraystretch}{1.3}
\caption{5-Fold CV Results – Classification}
\begin{tabular}{|1|c|c|c|p{4cm}|}
\hline
Model & Metric & Fold Avg (No-PCA) & Fold Avg (With-PCA) & Observation \\
\hline
Logistic Regression & Accuracy & 0.7340 & 0.7420 & Linear model benefits noticeably from PCA \\
Logistic Regression & F1-score & 0.7315 & 0.7403 & PCA removes redundant features improving F1 \\
SVM & Accuracy & 0.7430 & 0.7300 & SVM with RBF can handle high-dim, PCA may help speed \\
SVM & F1-score & 0.7367 & 0.7221 & Kernel SVM partially benefits from dimensionality reduction \\
\hline
\end{tabular}
\end{table}

%==================================================

\section{Analysis Questions}

\begin{itemize}
\item \textbf{Did PCA improve Linear Regression performance?}
    → MSE No PCA=26.2338  |  MSE With PCA=26.2251
    → YES: PCA reduced MSE.
      Linear Regression benefits from PCA when features are multicollinear,
      as PCA decorrelates them and removes noise dimensions.
 
\item \textbf{Did PCA significantly affect Random Forest?}
    → R² change = 0.0015 (No PCA=0.7225  |  With PCA=0.7241)
    → Change is minimal.
      Tree-based models use feature splits and are less sensitive to
      multicollinearity; PCA may even hurt by creating abstract components.

 
\item \textbf{Which classification model benefited more from PCA?}
    → Logistic Regression gain: +0.80% accuracy
    → SVM gain               : -1.30% accuracy
    → Logistic Regression benefited more from PCA.
 
\item \textbf{Did PCA reduce variance across folds?}
    → LR MSE std: No PCA=1.9532  |  With PCA=2.1705
    → NO: PCA did not stabilise fold performance.
      By removing noise dimensions, PCA often reduces variance in CV resu
 
\item \textbf{Why do linear models benefit more from PCA than ensemble models?}
    → Linear models assume feature independence and suffer from multicollinearity.
      PCA orthogonalises features (zero correlation), which directly helps.
      Ensemble (tree) models split one feature at a time, making them immune
      to multicollinearity. PCA's loss of interpretability (abstract PCs)
      doesn't compensate for any gain in tree-based methods.

\end{itemize}
%==================================================
\section{Expected Learning Outcome}

\begin{itemize}

\item PCA reduces multicollinearity
\item PCA helps linear models more
\item Tree-based models may not benefit much
\item PCA reduces noise
\item Dimensionality reduction can improve generalization

\end{itemize}


\section{Conclusion}
\begin{itemize}
\item \textbf{When PCA IS useful ?}
  • Linear models (Linear Regression, Logistic Regression, SVM-linear)
  • High-dimensional data with multicollinearity
  • When training speed and memory are constraints
  • When noise reduction is a priority
 
\item \textbf{When PCA is NOT necessary?}
  • Tree/ensemble models (Random Forest, XGBoost, Decision Trees)
  • When feature interpretability is required
  • When dataset already has few, independent features
 
\item \textbf{Trade-off:}
  • PCA improves generalization and reduces overfitting
  • BUT principal components are combinations of original features —
    interpretability is lost (PC1 is not "age" or "score", it is a
    weighted mix of all features)
  • Choose PCA when performance > interpretability
 
\item \textbf{Key Takeaways:}
  • PCA reduces multicollinearity
  • PCA helps linear models more than tree-based models
  • PCA can reduce noise and improve generalization
  • Retaining 95% variance is a good default strategy
\end{itemize}

\end{document}